% DRAFT (2026-09-24): written by Claude from the slides and the lecture notes;
% not yet reviewed by the author.
\section{Introduction}\label{sec:introduction}

% ── slides L1 p.4–12 ──
Artificial intelligence (AI) is a large family of methods. Machine learning
(ML) is one member of the family, deep learning (DL) is a branch of machine
learning, and the large language models (LLMs) behind tools such as ChatGPT and
Copilot are a branch of deep learning. This course covers, in the lecturer's
words, ``all AI, but learning'': methods that answer questions broader than
those of machine learning and that run, mostly unseen, inside many AI-based
tools. Some learning may still appear along the way, and
\cref{sec:intro-hybrid} gives the reason for this choice.

% Grouping of the lectures of p.5 into the four blocks of p.4; planning sits
% with logical agents as in AIMA part III (p.6); the slides place multi-agent
% systems in no block.
After a lecture on the design of intelligent agents (\cref{sec:agents}), the
course runs in four blocks:
\begin{itemize}
  \item \textbf{Solving problems by search}: uninformed and informed search,
    local search, constraint satisfaction problems, search in games.
  \item \textbf{Evolutionary optimization and artificial life}: evolutionary
    optimization, novelty-based search and quality-diversity, artificial life.
  \item \textbf{Logical agents}: inference in propositional and first-order
    logic, planning.
  \item \textbf{Probabilistic agents}: Bayesian networks, hidden Markov models.
\end{itemize}
Multi-agent systems close the course.

\subsection{Symbolic AI}\label{sec:intro-symbolic}

% ── slides L1 p.19 ──
Symbolic AI was the dominant paradigm of the field from the 1950s to the
1990s. It represents the world at a high level, in a form that humans can read
and interpret, and its building block is the symbol.

\begin{definition}[Symbol]\label{def:symbol}
  A \textbf{symbol} is a meaningful pattern that can be manipulated.
\end{definition}

A symbol stands for something else, as the string \texttt{cat} stands for the
animal, and a program can operate on it: compare it, copy it, combine it with
other symbols into rules. Semiotics, the study of signs and of how they carry
meaning, names the two sides of this relation: the pattern is the
\textbf{signifier}, and what it stands for is the \textbf{signified}.

Where do the symbols of a machine come from? The lecture gives two answers.
Either humans provide them, and then the machine's symbols mean something only
because someone who already had symbols gave them their meaning: a
chicken-and-egg problem. Or the computer builds its own low-level
representations from raw data, which is sub-symbolic manipulation
(\cref{sec:intro-subsymbolic}). The lecturer relates this circularity to
Hofstadter's ``strange loops''.

\subsection{The symbol grounding problem}\label{sec:intro-grounding}

% ── slides L1 p.20 ──
If a symbol means something only to the humans who supplied it, the system that
manipulates the symbol does not know what it manipulates. Harnad named this
difficulty.

\begin{definition}[Symbol grounding problem]\label{def:grounding}
  The \textbf{symbol grounding problem} asks how a symbol is connected to the
  real-world object it stands for: what makes the connection, who makes it,
  and how. In Harnad's words: ``How can the semantic interpretation of a
  formal symbol system be made intrinsic to the system, rather than just
  parasitic on the meanings in our heads?''\footnote{S.~Harnad, ``The symbol
  grounding problem'', \emph{Physica D} 42 (1990), pp.~335--346.}
\end{definition}

A symbol represents something outside the system, so the answer requires what
Hofstadter calls ``jumping out of the system''. Formal logic meets the same
limit in Gödel's incompleteness theorem: every consistent, effectively
axiomatized formal system that can express elementary arithmetic contains a
true statement that it cannot prove, so establishing that statement requires a
step outside the system.

\subsection{The physical symbol system hypothesis}\label{sec:intro-pssh}

% ── slides L1 p.21 ──
Suppose that symbols can be grounded. Is manipulating them then enough for
intelligence? Newell and Simon answered yes, in the form of a hypothesis about
a particular kind of system.

\begin{definition}[Physical symbol system]\label{def:pss}
  A \textbf{physical symbol system} is a system that manipulates symbols to
  produce other symbols. It is \emph{physical} because the manipulation is a
  physical process, which consumes energy.
\end{definition}

In the lecture's terms, humans are physical symbol systems, and so are digital
computers, since both manipulate symbols. The \textbf{physical symbol system
hypothesis} (PSSH) states:
\begin{quote}
  A physical symbol system has the necessary and sufficient means for general
  intelligent action.\footnote{A.~Newell and H.~A.~Simon, ``Computer science as
  empirical inquiry: symbols and search'', \emph{Communications of the ACM}
  19(3) (1976), pp.~113--126.}
\end{quote}
% The slide drops "general"; the wording and both glosses below follow the
% 1976 paper.
Newell and Simon explain both halves. \emph{Necessary} means that ``any system
that exhibits general intelligence will prove upon analysis to be a physical
symbol system''. \emph{Sufficient} means that ``any physical symbol system of
sufficient size can be organized further to exhibit general intelligence''.
The claim is empirical, which is why it is a hypothesis and not a theorem:
evidence can support it or refute it. If it holds, a digital computer, suitably
organized, can act with general intelligence.

\subsection{Intelligence is not consciousness}\label{sec:intro-consciousness}

% ── slides L1 p.22 ──
A system that acts intelligently need not be conscious: intelligence and
consciousness are different properties. Consciousness is usually tied to
subjective experience. In the sense of Nagel's essay ``What is it like to be a
bat?'',\footnote{T.~Nagel, ``What is it like to be a bat?'', \emph{The
Philosophical Review} 83(4) (1974).} a being is conscious if there is something
it is like to be that being. Subjective experience is hard to detect from
outside: dissecting a system and studying its behavior may not be enough. A
\textbf{zombie}, which behaves exactly like a person but has no experience, would
pass every behavioral test. Searle's Chinese room makes a related point about
understanding: a person who follows written rules to manipulate Chinese
characters can answer questions in Chinese without understanding any
Chinese.\footnote{J.~R.~Searle, ``Minds, brains, and programs'',
\emph{Behavioral and Brain Sciences} 3(3) (1980).}

If consciousness exists, materialism holds that it is a by-product of the body,
located ``somewhere'' in it; dualism instead treats the mind as a substance
different from the body. Hofstadter argues that what matters is the symbolic
pattern and not its physical substrate: silicon chips organized in an
equivalent way could support consciousness just as neurons do. The lecture
leaves the question open.

\subsection{Sub-symbolic representations}\label{sec:intro-subsymbolic}

% ── slides L1 p.24–25 (images only; the reading follows the slide titles) ──
Symbolic descriptions break down when the world varies more than a set of rules
can anticipate. The alternative is to let the machine build its own
representation.

\begin{definition}[Sub-symbolic representation]\label{def:subsymbolic}
  A \textbf{sub-symbolic representation} encodes a concept as a numeric pattern
  spread over many low-level units, such as the weights of a neural network,
  instead of as a single symbol that a human can read.
\end{definition}

\begin{example}[Recognizing a cat]\label{ex:cat}
  This example contrasts the two approaches. A symbolic system describes a cat
  through symbols for its parts, such as two eyes, whiskers and pointed ears,
  and through rules that combine them. Real photographs of cats vary in pose,
  breed, lighting and framing, and each new variation breaks some rule. A
  sub-symbolic system learns its representation from many example images and
  absorbs the variation, at the price of a representation that no human can
  read.
\end{example}

\subsection{Models of the world}\label{sec:intro-models}

% ── slides L1 p.26 (images only: a map of Pisa; house features pointing to a
% selling price; a path through a labyrinth) ──
Whatever its representation, an AI system reasons about a \textbf{model of the
world}: a representation of the part of reality that the system works on. The
lecture distinguishes three kinds of model, by the question that each answers.
\begin{itemize}
  \item \textbf{Descriptive}: how is the world? A map of Pisa describes where
    things are.
  \item \textbf{Predictive}: what value will an unobserved quantity take? A
    model estimates the selling price of a house from its floor area, its
    number of rooms and its energy efficiency. Machine learning builds models
    of this kind.
  \item \textbf{Decisional}: what should be done to reach a goal? A path through
    a labyrinth answers it; the search methods of the course answer questions
    of this kind.
\end{itemize}

\subsection{Hybrid systems}\label{sec:intro-hybrid}

% ── slides L1 p.27–28; p.29–31 (AI scientist, agentic web, evolutionary
% computation) are outside the agreed scope ──
Machine learning maps inputs $X$ to outputs $Y$. Building the map (training)
requires huge amounts of data, and the more complex the mapping, the more data
it requires. Using the map (inference) requires expensive memory, time and
dedicated accelerators such as GPUs and TPUs. Deep learning requires even more
data and computation. It still works very well in many applications, and the
lecture states the course's position explicitly: ``This course is not saying
`don't use deep learning'. It's saying `don't use deep learning only'.'' A
\textbf{hybrid system} combines a learned model with components that are not
learned.

\begin{example}[Retrieval-augmented generation]\label{ex:rag}
  Retrieval-augmented generation (RAG) is a worked case of a hybrid system. A
  \textbf{retriever} fetches the documents relevant to a question from an
  external memory, such as Wikipedia, and a \textbf{generator}, a language model,
  writes the answer from the question and those documents. We can expand the
  memory to a new domain, or update it with the news, without retraining the
  model, and each answer can point to the document it comes from, which makes
  it easier to interpret.
\end{example}
